% ── SECTION 2: BACKGROUND (~800 words) ───────────────────────────────────────

\section{Background}
\label{sec:background}

\subsection{Bell's Theorem and the End of Local Realism}

Bell's theorem~\cite{Bell1964} proves that no theory of local hidden variables
can reproduce all statistical predictions of quantum mechanics.
The argument is straightforward: if quantum particles possessed definite
pre-existing values for every observable~---~the hidden variable
hypothesis~---~those values would have to satisfy a set of linear inequalities
(the Bell inequalities) on the correlations between distant measurement outcomes.
Quantum mechanics predicts violations of these inequalities for entangled states.

The testable form due to Clauser, Horne, Shimony, and Holt~\cite{CHSH1969},
$|\mathcal{S}| = |E(a,b) - E(a,b') + E(a',b) + E(a',b')| \leq 2$,
was violated experimentally by Aspect and colleagues~\cite{Aspect1982}
and subsequently confirmed with all major loopholes closed by Hensen
\emph{et al.}\ and by Giustina \emph{et al.}\ in 2015, culminating in the
Nobel Prize-cited experiments of Zeilinger and colleagues~\cite{Zeilinger2023Nobel}.
The maximum quantum violation, $|\mathcal{S}| = 2\sqrt{2}$ (the Tsirelson bound),
is achieved by maximally entangled two-qubit states.

The physical implication is unambiguous: quantum particles do not possess
determinate properties prior to and independently of measurement.
Properties are not intrinsic to systems; they are relational~---~they emerge
through interaction.
This result~---~local hidden variable theories are false~---~is among the
most precisely confirmed facts in physics.
The \RCT~takes this conclusion as its foundational premise.

\subsection{Relational Quantum Mechanics}

Rovelli's relational quantum mechanics (RQM)~\cite{Rovelli1996,Rovelli2021}
radicalizes the implication of Bell's theorem at the level of quantum states
themselves.
In RQM, the quantum state $\ket{\psi}_{S/O}$ of system $S$ is defined only
relative to observer $O$~---~never absolutely.
A second observer $O'$, not yet correlated with the measurement outcome,
consistently assigns a different state to~$S$.
Neither description is more fundamental; both are valid relative to their
respective observational contexts.

The measurement event, in RQM, is the moment when a relational interaction
actualizes a definite outcome: it is not a collapse of an objective wave
function, but the establishment of a definite fact \emph{relative to} the
interacting system.
This account dissolves the measurement problem without invoking hidden variables,
many worlds, or epistemic fictions~---~it does so by taking the relational
character of quantum facts seriously as a feature of reality, not a limitation
of knowledge.

The critical lacuna in RQM, as its author acknowledges~\cite{Rovelli2021},
is the absence of a formal mathematical structure that makes relational
affinity~---~the degree of coherent coupling between systems~---~a computable,
dynamical quantity with experimentally testable consequences.
The \RCT~provides this structure.

\subsection{Decoherence and the Quantum-to-Classical Transition}

Environmental decoherence~\cite{Zurek2003} describes how quantum superpositions
are suppressed through entanglement with environmental degrees of freedom.
For system $S$ in a pure state $\ket{\psi_S}$ coupled to environment $E$,
the joint evolution produces an entangled state $\ket{\psi_{SE}}$.
The reduced density matrix of $S$,
\begin{equation}
  \rho_S = \mathrm{Tr}_E\bigl[\ket{\psi_{SE}}\bra{\psi_{SE}}\bigr],
\end{equation}
transitions from a pure state (with nonzero off-diagonal elements) to an
effectively mixed state (with suppressed off-diagonal elements) as
$S$-$E$ entanglement grows.
This suppression of coherence is not collapse~---~the full state
$\ket{\psi_{SE}}$ remains pure~---~but the \emph{relational coherence} of
$S$ with respect to any local observer is destroyed.

Zurek's einselection mechanism~\cite{Zurek2003} identifies a preferred
``pointer basis'' selected by the system-environment interaction Hamiltonian,
into which superpositions decohere.
The \RCT~reframes decoherence as the decay of the Affinity Tensor~$\alpha$
from its maximal value (full relational participation) to zero
(classical separability), providing a single scalar dynamical variable
that tracks the transition.

\subsection{Entanglement Measures}

Several measures quantify bipartite entanglement for two-qubit states.
The von Neumann entropy of entanglement,
\begin{equation}
  S(\rho_S) = -\mathrm{Tr}\bigl[\rho_S \log_2 \rho_S\bigr],
  \label{eq:vN-entropy}
\end{equation}
measures entanglement for pure bipartite states~\cite{vonNeumann1932,NielsenChuang2000}.
For mixed states, the concurrence~\cite{Wootters1998}
\begin{equation}
  \mathcal{C}(\rho) = \max\bigl(0,\, \lambda_1 - \lambda_2 - \lambda_3 - \lambda_4\bigr),
  \label{eq:concurrence}
\end{equation}
where $\lambda_i$ are the square roots of the eigenvalues (in decreasing order)
of $\rho(\sigma_y \otimes \sigma_y)\rho^*(\sigma_y \otimes \sigma_y)$,
provides an entanglement monotone defined on the full state space.

The Affinity Tensor~$\alpha(S,E)$ introduced in Section~\ref{sec:rct} is
related to these measures but defined for a different purpose.
While $S(\rho)$ and $\mathcal{C}(\rho)$ quantify entanglement at an instant,
$\alpha(t)$ is constructed to track the \emph{dynamics} of relational coherence
under environmental coupling~---~making it operationally suited to characterizing
gate performance in realistic noisy devices.
Section~\ref{sec:rct} establishes the precise relationship between $\bar\alpha$,
$\mathcal{C}$, and the off-diagonal structure of the relational density matrix.
